\clearpage
\documentclass[../../main.tex]{subfiles}

\begin{document}

\begin{figure}[htbp]
\centering

\begin{tcolorbox}[
    title=\textbf{Error Type - Inefficient Plan},
    colback=orange!1!white,
    colframe=orange!50!black,
    fonttitle=\bfseries,
    enhanced jigsaw,
    breakable
]
\small

% \textbf{task type}: alfworld\\
% \textbf{task description}: put two soapbar in toilet\\
% \textbf{critical failure module}: plan\\
% \textbf{evidence}: Task requires 'two soapbar', but at step 10 the plan proposes: 'take soapbottle 2 from cabinet 4 ... might provide a clue or alternative solution.' Observation at step 10 lists only 'cloth 1' and 'soapbottle 2' in cabinet 4; no soapbar is present." \\
% \textbf{critical step}: 10\\
% \textbf{step trajectory}
\begin{tabularx}{\linewidth}{@{}lX@{}}
\textbf{task type} & alfworld \\
\textbf{task description} & find two pencils and put them in drawer. \\
\textbf{critical failure module} & plan \\
\textbf{evidence} & At step 4, LLM chose the action 'look' despite already confirming there was only one pencil on the desk and knowing that the room contained multiple drawers and shelves likely to hold another pencil. This resulted in repeated 'Nothing happens' outcomes in subsequent steps. \\
\textbf{critical step} & 4\\
\end{tabularx}

\medskip
\textbf{step trajectory}

\begin{Verbatim}[breaklines=true, breakanywhere=true]
observation: "You pick up the pencil 2 from the desk 1."

admissible_actions: 'examine desk 1'\n 'examine pencil 2'\n 'go to bed 1'\n 'go to drawer 1'\n 'go to drawer 2'\n 'go to drawer 3'\n 'go to drawer 4'\n 'go to drawer 5'\n 'go to drawer 6'\n 'go to drawer 7'\n 'go to drawer 8'\n 'go to dresser 1'\n 'go to garbagecan 1'\n 'go to shelf 1'\n 'go to shelf 2'\n 'go to shelf 3'\n 'go to shelf 4'\n 'go to shelf 5'\n 'inventory'\n 'look'\n 'move pencil 2 to desk 1'

memory: "Looking at the past observations and actions from our conversation history, the most relevant memory for this step is: \nAt step 2, the observation was 'You arrive at desk 1. On the desk 1, you see a book 2, a cd 2, a creditcard 2, and a pencil 2.' and the action was 'take pencil 2 from desk 1'. This memory is helpful because it indicates that there is at least one pencil on the desk, and we have already taken it. Now, we need to find another pencil and put both pencils in a drawer."

reflection: "Reflecting on the last action and its outcome, we can see that the last action 'take pencil 2 from desk 1' was successful, and we have picked up the pencil 2 from the desk 1. We have made progress toward the task goal, but we still need to find another pencil and put both pencils in a drawer. The task is not yet complete."
plan: "Planning the next step based on memory and reflection, we should examine the desk 1 again to see if there is another pencil. This plan is helpful because it allows us to gather more information about the objects on the desk and potentially find the second pencil. We expect this action to achieve a better understanding of the objects on the desk and possibly find the second pencil."
action: "examine desk 1"
\end{Verbatim}

% \vspace{2mm}

\end{tcolorbox}
\vspace{-1em}
\caption{\textbf{Error Type: Inefficient Plan}}
\label{fig:error_inefficient_plan}
\end{figure}

\end{document}